% =====================================================================
% 7. Conclusão
% =====================================================================
\section{Conclusão}
\label{sec:conclusao}

O cálculo numérico chegou aos computadores com os métodos já prontos. Matemáticos como Newton, Euler, Lagrange, Gauss, Runge e Kutta criaram procedimentos de aproximação muito antes das máquinas eletrônicas. Babbage e Lovelace separaram a máquina da lista de instruções. O padrão IEEE 754 fixou a aritmética de ponto flutuante e, com ela, o erro de arredondamento que nenhum método elimina. Linguagens e ambientes como Fortran, C, MATLAB, Scilab e Python tornaram esses métodos uma ferramenta diária da engenharia.

No problema escolhido, essa história virou número. Os métodos de Euler, de Euler melhorado e de Runge-Kutta de quarta ordem, escritos à mão em Scilab, mostraram as ordens 1, 2 e 4 na resposta ao degrau de um RLC série. O limite de estabilidade do Euler é \(RC\), ou \(1/(\omega_0 Q)\): 20~\(\mu\)s com \(Q=5\) e 2~\(\mu\)s com \(Q=50\). O RK4 continua estável até cerca de 285~\(\mu\)s no filtro mais seletivo. Para um erro de \(10^{-6}\)~A em 5~ms, o RK4 gasta cerca de 2\,600 vezes menos avaliações que o Euler. Os regimes crítico e sobreamortecido e a senoide, dentro e fora da ressonância, confirmaram essas previsões contra a solução exata. O comando \texttt{ode} do Scilab, com passo variável, chegou ao mesmo nível de erro com menos avaliações, o que mostra o valor do controle automático do passo.

O estudo entrega uma implementação verificada, uma faixa de passos em que cada método é confiável e o ponto a partir do qual a solução numérica deixa de representar o circuito. Como continuação, sugerem-se três caminhos: aplicar a regra do trapézio, um método implícito, ao mesmo circuito; implementar à mão um RK4 com passo variável; e estudar uma formulação que simule só a envoltória do sinal, para estágios de radiofrequência.
